

\begin{abstract}
本文针对零售销售额预测与数据分析问题，提出了基于描述性统计、多维方差分析与Holt-Winters时间序列预测的系统解决方案。

\textbf{针对问题一：}本问是描述性统计与可视化问题，需对零售数据集进行统计分析和多维度可视化。首先，对数据集加载并处理缺失值。接着，计算销售额核心统计量，均值230.77美元、中位数54.49美元、标准差626.65美元，总销售额226.15万美元、总订单数4921个、有效客户数793个。然后，按日、月、年三级聚合销售额绘制趋势图，按区域、产品大类、客户细分等多维度分析分布特征。结果表明，销售额呈高度右偏分布，偏度12.98，年度销售额从2022年47.99万美元增长至2025年72.21万美元，\textbf{West区域和Technology品类}是主要贡献来源。

\textbf{针对问题二：}本问是影响因素相关性分析问题，需量化各维度与销售额的关联程度。首先，对分类变量采用ANOVA方差分析，计算F统计量和$\eta^2$效应量。接着，计算月度数值特征间的Spearman秩相关系数。然后绘制热力图、箱线图和对比柱状图。结果表明，\textbf{产品子类别（$\eta^2$=0.199）和产品大类（$\eta^2$=0.051）}是影响销售额最显著的因素；月度层面，\textbf{订单数和客户数与销售额高度相关（$r_s>0.9$）}。

\textbf{针对问题三：}本问是销售额预测建模问题，需构建预测模型并评估。首先按月聚合构建月度销售额序列，提取时间、滞后和滚动窗口特征。接着，对比季节性朴素预测、Holt-Winters指数平滑、随机森林和XGBoost四种模型。以2022-2024年为训练集、2025年为测试集，采用MSE、MAE、MAPE评估。结果表明，\textbf{Holt-Winters模型最优，MAPE为23.56\%，$R^2$为0.7215}，有效捕捉了年度季节性和上升趋势。

\textbf{针对问题四：}本问是未来预测与策略建议问题。基于最优Holt-Winters模型进行2026年12个月滚动预测，预计全年总销售额89.94万美元，同比增长24.6\%，Q4达全年峰值。基于预测与影响因素分析，从产品布局、区域营销、客户运营、物流履约、库存管理五维度提出了可落地的经营策略。

综上，本文为零售销售额预测问题提供了从数据探索到预测建模再到业务应用的完整解决方案。

\keywords{零售销售额预测；Holt-Winters指数平滑；方差分析；时间序列分析；多维度数据分析}
\label{abstract:end}
\end{abstract}
